\chapter{How Deep Reinforcement Learning Works}

This chapter assumes nothing. If you already know actor--critic methods, skip to
\S\ref{sec:sacdetail}, where SAC is described as this project uses it.

\section{Learning by consequence}

Most machine learning is \emph{supervised}: you are shown thousands of examples
with the right answer attached, and you learn to reproduce them. That is
unavailable here --- nobody can supply the correct muscle forces, because that is
the very thing we cannot measure.

Reinforcement learning removes the need for correct answers. Instead:

\begin{enumerate}[leftmargin=1.5em, itemsep=2pt]
\item The learner observes the current situation.
\item It chooses an action.
\item The world changes, and it receives a number --- the \emph{reward} --- saying
how good that was.
\item Repeat, millions of times, adjusting so that high-reward behaviour becomes
more likely.
\end{enumerate}

\begin{plain}
It is learning the way you learn to ride a bicycle: nobody tells you the correct
handlebar angle. You try, you wobble, you notice which corrections stop the
wobbling, and you do more of those. The wobbling is the reward signal.
\end{plain}

\section{The vocabulary, applied to this arm}

\begin{longtable}{@{} L{3.0cm} L{5.2cm} L{7.4cm} @{}}
\toprule
\textbf{Term} & \textbf{General meaning} & \textbf{In this project} \\
\midrule
\endhead
Agent & the thing being trained & a neural network, 393{,}750 numbers \\
Environment & everything it acts on & the simulated arm and a target point \\
State & the true situation & joint angles, speeds, muscle states \\
Observation & what the agent is told & 38 numbers (\S\ref{sec:obs}) \\
Action & what it outputs & 9 numbers in $[0,1]$, one per muscle \\
Reward & the score after each action & a formula, seven terms (\S\ref{sec:reward}) \\
Episode & one attempt & one 2-second reach, 100 decisions \\
Policy & the learned rule & observation $\rightarrow$ action; this \emph{is} the trained controller \\
Return & total reward over an episode & what the agent actually maximises \\
Value & expected return from a state & the critic's estimate (\S\ref{sec:critic}) \\
\bottomrule
\end{longtable}

\begin{why}
Note the distinction between \emph{state} and \emph{observation}. The state is
what is true; the observation is what the agent is permitted to see. Here they
differ deliberately: the agent sees muscle lengths, forces and activations, joint
angles and speeds, and the direction to the target --- but not, for instance, the
randomised body masses used during training. It must cope without knowing them,
which is what makes the controller robust rather than brittle.
\end{why}

\subsection{The reward is a scoring rule, not an instruction}

This distinction causes more trouble than any other idea in the subject, so it is
worth stating sharply.

The agent is \textbf{never told how to move}. It is told only how good the
outcome was. Any behaviour that scores highly will therefore be found, including
behaviour the designer never imagined and would not endorse.

\begin{trap}
In this project the reward was, for a period, maximised by \textbf{deliberately
destroying the simulated arm}. Every per-step reward was zero or negative, and a
numerical blow-up ended the episode with no penalty. In this framework an episode
that \emph{ends} accrues nothing further --- its remaining value is exactly zero
--- while surviving the full 100 steps accumulated about $-20$. Quitting was
therefore worth more than participating.

The agent found this. Measured over 50 evaluation episodes: \textbf{100\,\% ended
in a deliberate blow-up at step 7 of 100.} Not one ran to completion.

The algorithm was not broken. It did exactly what it was asked. The request was
wrong. Full account in \S\ref{sec:defect3}.
\end{trap}

\section{Why a ``value'' is needed, and what the discount factor does}
\label{sec:critic}

An action's consequences are not immediate. Pulling a muscle now changes where
the hand is half a second later. So the agent cannot judge an action by the
reward that immediately follows it --- it must judge by the reward that
\emph{eventually} follows.

That total is the \textbf{return}:
\[
G_t \;=\; r_t \;+\; \gamma\, r_{t+1} \;+\; \gamma^2 r_{t+2} \;+\; \gamma^3 r_{t+3} + \cdots
\]

The number $\gamma$ (gamma), between 0 and 1, is the \textbf{discount factor}. It
shrinks rewards that arrive later.

\begin{plain}
$\gamma$ answers: ``how far into the future should I care?''

At $\gamma = 0$ the agent is purely greedy --- only the next reward matters. At
$\gamma$ close to 1 it weighs events far ahead almost as heavily as now.

A useful rule of thumb: the agent effectively plans about $1/(1-\gamma)$ steps
ahead. At $\gamma = 0.99$ that is 100 steps. At $\gamma = 0.957$ it is 23.
\end{plain}

\begin{why}
Remember that rule of thumb. \textbf{The discount factor turned out to be the
single most consequential setting in this entire project} --- worth more than the
choice of algorithm --- and the reason is exactly this horizon arithmetic. The
default value made the agent plan over 100 steps when the reach it was learning
completes in about 25. Chapter~\ref{ch:results}, \S\ref{sec:gammaresult}.
\end{why}

Because the future is unknown, the agent learns to \emph{estimate} the return.
That estimator is the \textbf{critic} (or value function): given a situation and
an action, it predicts the total reward that will follow. The critic is what
makes learning possible from delayed consequences --- and, as we will see, a
critic fed noisy targets is what makes learning fail.

\section{Actor and critic}

Modern methods train two networks together:

\begin{itemize}[leftmargin=1.4em, itemsep=3pt]
\item The \textbf{actor} (the policy) chooses actions. This is the thing we
ultimately keep and deploy.
\item The \textbf{critic} estimates how good those choices are.
\end{itemize}

They bootstrap each other: the critic learns to predict returns from experience,
and the actor adjusts to prefer actions the critic rates highly.

\begin{plain}
A coach and an athlete. The athlete performs; the coach judges. The athlete
changes technique to earn better judgements; the coach refines their judgement by
watching what actually happens. Both improve together --- and if the coach is
unreliable, the athlete trains toward the wrong thing. That failure mode is not
hypothetical here; it is what the discount factor caused.
\end{plain}

\section{Off-policy learning and the replay buffer}

An \textbf{off-policy} method stores every experience in a large memory --- the
\emph{replay buffer} --- and re-learns from it repeatedly. An \textbf{on-policy}
method uses each batch of experience once and discards it.

\begin{plain}
Off-policy keeps a notebook of every attempt ever made and re-reads it
constantly. On-policy makes a batch of attempts, learns from them, throws the
notebook away, and starts fresh.

Re-reading is far more efficient when attempts are expensive, which is why three
of our four algorithms do it.
\end{plain}

The \textbf{replay ratio} --- gradient updates performed per environment step ---
matters enormously and is easy to get wrong.

\begin{trap}
This project ran with four environments in parallel and one gradient update per
rollout iteration. Because one iteration over four environments collects
\emph{four} transitions, the actual ratio was $0.233$ updates per step rather than
the standard $\approx 1$. Every experiment was doing roughly a quarter of its
intended training, invisibly. Measured directly in \S\ref{sec:defect5}.
\end{trap}

\section{The four algorithms}
\label{sec:sacdetail}

\subsection{DDPG --- Deep Deterministic Policy Gradient}

The oldest of the four. One actor, one critic. The actor outputs a single
definite action; exploration is added by injecting random noise into it.

Its known weakness is \emph{overestimation}: the critic systematically rates
actions better than they are, the actor chases those inflated ratings, and
training becomes unstable.

\subsection{TD3 --- Twin Delayed DDPG}

DDPG with three corrections. It trains \textbf{two} critics and always believes
the more pessimistic one, which counters overestimation. It updates the actor
less often than the critics (``delayed''), and it smooths the critic's target by
adding small noise, so the actor cannot exploit a sharp spike in the critic's
estimate.

\subsection{SAC --- Soft Actor-Critic}

Also uses twin critics, but differs in one important way: its policy is
\textbf{stochastic}. It outputs a probability distribution over actions rather
than a single action, and it is explicitly rewarded for keeping that distribution
spread out. The quantity being maximised is
\[
\sum_t \Big( r_t \;+\; \alpha\, \mathcal{H}[\pi(\cdot \mid s_t)] \Big),
\]
where $\mathcal{H}$ is entropy --- a measure of randomness --- and $\alpha$ its
weight.

\begin{plain}
Entropy is a measure of how undecided something is. A policy that always does
exactly the same thing has zero entropy. One that picks randomly has high
entropy.

SAC pays the agent a bonus for staying undecided. The purpose is to stop it
committing early to a mediocre habit before exploring alternatives.

The weight $\alpha$ is tuned automatically to hit a chosen \textbf{target
entropy}. The library default sets that target to minus the number of controls
--- here $-9$, because there are nine muscles.
\end{plain}

\begin{why}
That default is worth pausing on. It scales with the \emph{number of actions},
not with anything about the task. A nine-muscle arm gets more forced randomness
than a two-joint one, regardless of whether the task needs precision.

For a long time this project believed that default was the cause of its
inability to hold position --- the mechanism fits perfectly, and every one of the
best hyperparameter trials agreed on a much lower value. \textbf{A direct
ablation later proved that belief wrong.} The story is told in
\S\ref{sec:entropyrefuted}, because being wrong for a good reason is instructive.
\end{why}

\subsection{PPO --- Proximal Policy Optimisation}

Structurally different: on-policy. It collects a batch of experience, improves
the policy while limiting how far it may change in one step (the ``proximal''
constraint, which prevents destructive updates), then discards the batch.

Simple and robust, but much less sample-efficient --- it needs far more
interaction for the same skill, which matters when each interaction requires
simulating physics.

\subsection{Summary of the differences that matter here}

\begin{center}
\begin{tabular}{@{} L{2.2cm} C{2.0cm} C{2.2cm} C{2.6cm} L{4.6cm} @{}}
\toprule
& \textbf{Critics} & \textbf{Policy} & \textbf{Reuses data} & \textbf{Exploration by} \\
\midrule
DDPG & 1 & deterministic & yes & added noise \\
TD3  & 2 (min) & deterministic & yes & added noise \\
SAC  & 2 (min) & stochastic & yes & entropy bonus \\
PPO  & 1 (value) & stochastic & no & policy's own randomness \\
\bottomrule
\end{tabular}
\end{center}

\begin{why}
Only SAC has a \emph{target entropy}, because only SAC has an entropy bonus to
tune. This asymmetry matters for the algorithm comparison: there is no single
setting one can hold equal across all four, which is part of why fair comparison
between reinforcement learning algorithms is harder than it looks
(\S\ref{sec:fairness}).
\end{why}
